\subsection*{File 05H: Historical Cosine Similarity Defence with Warm-Up Phase}

File 05H extends the poisoning defence investigation by replacing the Median/MAD defence from File 05G with a historical cosine-similarity-based anomaly detection mechanism. The objective of this experiment was to determine whether directional similarity between client updates across communication rounds could provide a more adaptive and privacy-preserving poisoning defence for both Federated Learning (FL) and Split Federated Learning (SFL).

The experiment continues using:
\begin{itemize}
    \item the MNIST dataset,
    \item the extreme one-label-per-client non-IID configuration,
    \item the three-layer MLP architecture,
    \item and the same Scale Gradient Attack (SGA) poisoning attack introduced in File 05F.
\end{itemize}

The model architecture remains:

\[
784 \rightarrow 256 \rightarrow 128 \rightarrow 10
\]

with the SFL split defined as:

\[
\text{Client side: } 784 \rightarrow 256
\]

\[
\text{Server side: } 256 \rightarrow 128 \rightarrow 10
\]

Unlike File 05G, which used update magnitudes and Median/MAD robust statistics, File 05H attempts to detect poisoning based on directional inconsistency of client updates over time.

\subsubsection*{Motivation for Historical Cosine Defence}

The motivation behind cosine-similarity-based defence was that poisoning attacks often change the direction of optimisation updates, not only their magnitude.

In normal training:
\begin{itemize}
    \item client updates should gradually follow relatively consistent optimisation directions,
    \item especially when local data distributions remain stable across rounds.
\end{itemize}

A malicious client performing SGA poisoning introduces highly abnormal directional changes, even if update magnitudes are clipped or controlled.

Therefore, instead of measuring only update size, File 05H measures whether the direction of a client's update remains consistent with its own historical behaviour.

\subsubsection*{Historical Warm-Up Phase}

One major challenge in cosine-based defences is that early training rounds naturally contain unstable gradients.

During initial optimisation:
\begin{itemize}
    \item honest client updates can differ significantly,
    \item particularly under extreme non-IID conditions.
\end{itemize}

To reduce false positives, a warm-up phase was introduced.

During the warm-up period:
\begin{itemize}
    \item no clients are flagged,
    \item and historical memory vectors are gradually accumulated.
\end{itemize}

This allows the defence to first learn:
\begin{itemize}
    \item the normal optimisation behaviour of each client.
\end{itemize}

Without this warm-up phase:
\begin{itemize}
    \item the defence frequently classified honest clients as malicious during early rounds.
\end{itemize}

\subsubsection*{Mathematics of Cosine Similarity}

For each client update vector:

\[
u_k^{(t)}
\]

the cosine similarity with the client's historical reference vector:

\[
h_k^{(t-1)}
\]

is computed as:

\[
\cos(\theta_k)
=
\frac{
u_k^{(t)}
\cdot
h_k^{(t-1)}
}{
\|u_k^{(t)}\|
\,
\|h_k^{(t-1)}\|
}
\]

where:
\begin{itemize}
    \item values close to \(1\) indicate highly aligned updates,
    \item values near \(0\) indicate orthogonal behaviour,
    \item and negative values indicate opposite optimisation directions.
\end{itemize}

Clients with unusually low cosine similarity are considered suspicious.

\subsubsection*{Historical EMA Reference Vector}

The historical reference vector is updated using an Exponential Moving Average (EMA):

\[
h_k^{(t)}
=
\alpha
h_k^{(t-1)}
+
(1-\alpha)
u_k^{(t)}
\]

where:
\begin{itemize}
    \item \(h_k^{(t)}\) is the historical memory vector,
    \item \(u_k^{(t)}\) is the current update,
    \item and \(\alpha\) is the EMA smoothing coefficient.
\end{itemize}

This allows the defence to:
\begin{itemize}
    \item retain long-term optimisation behaviour,
    \item while still adapting gradually to changing gradients.
\end{itemize}

The historical EMA mechanism was substantially more stable than:
\begin{itemize}
    \item simple round-to-round cosine comparison,
    \item or direct global-reference cosine comparison.
\end{itemize}

\subsubsection*{Other Cosine Defences Attempted}

Several alternative cosine-based approaches were tested before settling on the historical warm-up implementation:

\begin{enumerate}
    \item \textbf{Direct round-to-round cosine similarity}
    \begin{itemize}
        \item compared each update only against the previous round,
        \item but became extremely unstable under non-IID distributions.
    \end{itemize}

    \item \textbf{Global-model cosine comparison}
    \begin{itemize}
        \item compared client updates against the global model update,
        \item but honest non-IID clients naturally diverged strongly from the global direction.
    \end{itemize}

    \item \textbf{Pure cosine thresholding without warm-up}
    \begin{itemize}
        \item produced excessive false positives during early rounds.
    \end{itemize}

    \item \textbf{Fixed cosine thresholds}
    \begin{itemize}
        \item failed because cosine distributions changed substantially throughout training.
    \end{itemize}
\end{enumerate}

The historical EMA warm-up version performed better because:
\begin{itemize}
    \item it personalised the reference behaviour for each client,
    \item and reduced early-round instability.
\end{itemize}

\subsubsection*{How Cosine Defence was Applied in FL}

In FL:
\begin{itemize}
    \item the attack affects all three model layers,
    \item therefore the cosine similarity is computed on the full flattened model update vector.
\end{itemize}

For each FL client:

\[
u_k^{FL}
=
\theta_k^{(t)}
-
\theta^{(t-1)}
\]

The cosine score is computed using the client's full update vector against its EMA historical memory.

Flagged clients are excluded from aggregation:

\[
\mathcal{K}_{keep}
=
\{
k \mid \cos(\theta_k) \geq \tau
\}
\]

The remaining honest clients are aggregated normally.

\subsubsection*{How Cosine Defence was Applied in SFL}

In SFL:
\begin{itemize}
    \item only the client-side layer is directly poisoned by the attacker.
\end{itemize}

Therefore:
\begin{itemize}
    \item cosine similarity is computed only on the client-side FL update portion.
\end{itemize}

This means the defence evaluates:
\begin{itemize}
    \item client-observable optimisation behaviour,
    \item without exposing raw data,
    \item activations,
    \item or labels.
\end{itemize}

Importantly:
\begin{itemize}
    \item the SFL server still does not directly access private client datasets.
\end{itemize}

Only optimisation statistics are exchanged.

The SFL server-side aggregation remains separate:

\[
\theta_s^{t+1}
=
\sum_{k=1}^{K}
\frac{n_k}{n}
\theta_{s,k}^{t+1}
\]

while the client-side FL aggregation uses cosine-filtered updates.

\subsubsection*{Results and Interpretation}

Compared to File 05F (no defence):
\begin{itemize}
    \item File 05H clearly improves SFL robustness.
\end{itemize}

The SFL model:
\begin{itemize}
    \item eventually recovers,
    \item and reaches approximately \(0.58-0.60\) accuracy.
\end{itemize}

This represents a major improvement over:
\begin{itemize}
    \item the undefended poisoning experiment,
    \item where optimisation became substantially more unstable.
\end{itemize}

However, the FL system still collapses catastrophically.

The FL accuracy falls toward:

\[
0.10
\]

which corresponds approximately to random guessing on MNIST.

\begin{figure}[H]
    \centering
    \includegraphics[width=0.85\linewidth]{figures/05H_accuracy.png}
    \caption{FL vs reconstructed SFL global test accuracy using the historical cosine defence with warm-up. SFL partially recovers, while FL still collapses under poisoning.}
\end{figure}

The FL loss also explodes dramatically:

\[
\text{loss} \rightarrow 70
\]

indicating:
\begin{itemize}
    \item severe optimisation divergence,
    \item and failure of the cosine defence to sufficiently isolate poisoned FL updates.
\end{itemize}

\begin{figure}[H]
    \centering
    \includegraphics[width=0.85\linewidth]{figures/05H_loss.png}
    \caption{FL vs reconstructed SFL global test loss under historical cosine defence. FL diverges catastrophically while SFL remains comparatively stable.}
\end{figure}

\subsubsection*{Why Cosine Defence Still Did Not Work Well}

Although cosine similarity improved SFL stability somewhat, the overall defence remained substantially weaker than the Median/MAD defence from File 05G.

Several reasons explain this:

\begin{enumerate}
    \item \textbf{Extreme non-IID behaviour naturally changes update directions}
    \begin{itemize}
        \item honest clients frequently appear directionally inconsistent.
    \end{itemize}

    \item \textbf{Cosine similarity ignores update magnitude}
    \begin{itemize}
        \item a poisoned update can still align directionally while having destructive magnitude.
    \end{itemize}

    \item \textbf{Small client count}
    \begin{itemize}
        \item only 10 clients were available,
        \item reducing statistical robustness.
    \end{itemize}

    \item \textbf{Deep model splitting}
    \begin{itemize}
        \item SFL client-side updates already contain reduced information compared to full FL updates.
    \end{itemize}

    \item \textbf{Historical drift}
    \begin{itemize}
        \item the EMA memory gradually adapts,
        \item meaning malicious updates can slowly contaminate historical references.
    \end{itemize}
\end{enumerate}

Overall, File 05H demonstrated that:
\begin{itemize}
    \item cosine-based directional defences alone are insufficient under strong poisoning attacks and extreme non-IID settings,
    \item although SFL naturally remains more robust than FL,
    \item and historical warm-up mechanisms significantly outperform naive cosine-threshold methods.
\end{itemize}
\begin{figure}[H]
    \centering
    \includegraphics[width=0.8\linewidth]{figures/historical_cosine_defence.png}
    \caption{Historical cosine similarity defence using EMA memory vectors.}
    \label{fig:historical_cosine}
\end{figure}
The experiment also provided strong motivation for the later combined defence approaches implemented in Files 05I and 05J, where:
\begin{itemize}
    \item cosine similarity was integrated with Median/MAD statistics,
    \item rather than being used independently.
\end{itemize}